\documentclass[12pt]{article}
\usepackage{amsmath,amssymb}

\begin{document}
\section*{{2020 AMC 10B Problems/Problem 10}}
Source: AoPS Wiki

\subsection*{Solution}
Notice that when the cone is created, the 2 shown radii when merged will become the slant height of the cone and the intact circumference of the circle will become the circumference of the base of the cone. We can calculate that the intact circumference of the circle is $8\pi\cdot\frac{3}{4}=6\pi$ . Since that is also equal to the circumference of the cone, the radius of the cone is $3$ . We also have that the slant height of the cone is $4$ . Therefore, we use the Pythagorean Theorem to calculate that the height of the cone is $\sqrt{4^2-3^2}=\sqrt7$ . The volume of the cone is $\frac{1}{3}\cdot\pi\cdot3^2\cdot\sqrt7=\boxed{\textbf{(C)}\ 3 \pi \sqrt7 }$ -PCChess

\end{document}
